\documentclass[11pt]{article}
\usepackage[spanish]{babel}
\usepackage[utf8]{inputenc}
\usepackage[T1]{fontenc}
\usepackage[a4paper,margin=2.5cm]{geometry}
\usepackage{enumitem}
\usepackage{hyperref}
\usepackage{lmodern}

\title{Preguntas de preparación para examen de candidatura}
\author{John Jairo Leal Gómez}
\date{\today}

\begin{document}

\maketitle

\section*{Uso sugerido}
Este documento reúne preguntas probables para preparar el examen de candidatura doctoral, organizadas por áreas de posible evaluación. La idea es usarlas como banco de preguntas para simulacros, respuestas cortas y respuestas ampliadas.

\section{Área ambiental}

\subsection*{Enfoque general}
\begin{itemize}[leftmargin=*]
    \item ¿Dónde está lo ambiental en la tesis y no solo lo técnico?
    \item ¿Cómo conocer las emisiones mejora realmente la gestión ambiental?
    \item ¿Por qué la medición de emisiones es relevante para sostenibilidad y no solo para monitoreo?
    \item ¿Qué consecuencias trae no conocer las emisiones o conocerlas de forma fragmentada?
    \item ¿Qué aporta esta tesis a la mitigación del cambio climático?
    \item ¿Cómo se conecta el proyecto con la gestión sostenible de residuos?
\end{itemize}

\subsection*{Gestión ambiental y alcance}
\begin{itemize}[leftmargin=*]
    \item ¿Por qué la tesis no se formuló directamente como un Sistema de Gestión Ambiental?
    \item ¿Qué relación tiene tu propuesta con un futuro SGA?
    \item ¿Qué parte de la gestión ambiental habilita tu sistema?
    \item ¿Cuál es el límite entre monitoreo ambiental y gestión ambiental en tu proyecto?
    \item ¿Cómo se pasa del dato ambiental a la toma de decisiones?
\end{itemize}

\subsection*{MRV, trazabilidad y sostenibilidad}
\begin{itemize}[leftmargin=*]
    \item ¿Por qué el enfoque de medición, reporte y validación es importante en este contexto?
    \item ¿Qué gana un sistema de bioconversión al tener trazabilidad de emisiones?
    \item ¿Cómo evita tu propuesta decisiones ciegas o reportes débiles?
    \item ¿Qué relación existe entre MRV y mercados de carbono?
    \item ¿Cómo evita el proyecto riesgos de greenwashing?
\end{itemize}

\subsection*{Riesgos, ética y tecnología}
\begin{itemize}[leftmargin=*]
    \item ¿Qué riesgos trae incorporar IA e IoT en sistemas de gestión de residuos?
    \item ¿Qué implicaciones éticas tiene el uso de datos en este proyecto?
    \item ¿Cómo abordas la opacidad algorítmica?
    \item ¿Qué significa soberanía de datos en el contexto del productor local?
    \item ¿Qué limitaciones ambientales y sociales puede tener una implementación a escala real?
\end{itemize}

\subsection*{Dimensión social y económica}
\begin{itemize}[leftmargin=*]
    \item ¿Cómo dialoga esta tesis con la dimensión social de los estudios ambientales?
    \item ¿Qué barreras económicas tiene escalar esta propuesta?
    \item ¿Qué retos de adopción enfrentaría el sistema en escenarios reales?
    \item ¿Qué oportunidades tendría para productores agroindustriales?
    \item ¿Cómo se conecta el proyecto con economía circular?
\end{itemize}

\section{Área técnica y tecnológica}

\subsection*{Arquitectura del sistema}
\begin{itemize}[leftmargin=*]
    \item ¿Por qué la arquitectura IoT propuesta es adecuada para este problema?
    \item ¿Cómo se integran sensor, microcontrolador, comunicación, backend, base de datos y dashboard?
    \item ¿Por qué escogiste esa arquitectura y no una solución más simple?
    \item ¿Cuál es el cuello de botella más probable del sistema?
    \item ¿Qué parte del sistema consideras más crítica para garantizar confiabilidad?
\end{itemize}

\subsection*{Sensores y adquisición de datos}
\begin{itemize}[leftmargin=*]
    \item ¿Por qué seleccionaste esos sensores para CO$_2$, CH$_4$ y variables ambientales?
    \item ¿Cómo vas a calibrar los sensores?
    \item ¿Contra qué referencia validarás las mediciones?
    \item ¿Qué significa en la práctica un error relativo menor al 5\%?
    \item ¿Cómo manejarás deriva, ruido, interferencia por humedad y pérdida de señal?
    \item ¿Qué variable es más crítica para la confiabilidad del sistema y por qué?
\end{itemize}

\subsection*{Modelado y sistema híbrido}
\begin{itemize}[leftmargin=*]
    \item ¿Por qué necesitas un enfoque híbrido y no solo IA o solo modelado matemático?
    \item ¿Qué entiendes exactamente por sistema híbrido en esta tesis?
    \item ¿Qué papel cumplen sensores, modelos matemáticos e IA dentro del sistema?
    \item ¿Qué parte del problema resolverá el modelo matemático y cuál la IA?
    \item ¿Cómo justificarías el uso de PINN frente a otros modelos más clásicos?
    \item ¿Cómo incorporarías la estocasticidad del proceso en el modelado?
\end{itemize}

\subsection*{Datos y validación}
\begin{itemize}[leftmargin=*]
    \item ¿Cómo construirás la base de datos de entrenamiento?
    \item ¿Qué harás si el volumen de datos resulta insuficiente?
    \item ¿Cómo separarás entrenamiento, validación y prueba?
    \item ¿Por qué elegiste métricas como MSE y \(R^2\)?
    \item ¿Cómo demostrarás que el modelo mejora la medición y no solo ajusta ruido?
    \item ¿Cómo distinguirás anomalías del sensor frente a anomalías reales del proceso?
\end{itemize}

\subsection*{Dashboard y trazabilidad}
\begin{itemize}[leftmargin=*]
    \item ¿Cuál es el valor científico del dashboard y no solo su valor visual?
    \item ¿Qué indicadores debe mostrar el dashboard para ser útil?
    \item ¿Cómo asegurarás integridad temporal y consistencia de datos?
    \item ¿Cómo versionarás calibraciones, cambios de sensor y cambios de modelo?
    \item ¿Cómo se conecta el dashboard con el componente MRV?
\end{itemize}

\subsection*{Viabilidad y alcance técnico}
\begin{itemize}[leftmargin=*]
    \item ¿Cuál es el producto mínimo funcional que esperas tener al final del doctorado?
    \item ¿Qué diferencia hay entre prototipo funcional, sistema validado y solución escalable?
    \item ¿Qué parte del proyecto ya está madura y cuál sigue en desarrollo?
    \item ¿Qué principal riesgo técnico podría comprometer resultados?
    \item ¿Cómo defenderías el proyecto si al final solo logras una implementación parcial?
\end{itemize}

\section{Área de mosca soldado negra y bioconversión}

\subsection*{Proceso biológico}
\begin{itemize}[leftmargin=*]
    \item ¿Cómo transforma \textit{Hermetia illucens} los residuos orgánicos?
    \item ¿Qué etapas del proceso biológico son más relevantes para tu investigación?
    \item ¿Por qué el metabolismo larval puede entenderse como fuente dinámica de gases?
    \item ¿Qué diferencia hay entre degradación del sustrato, crecimiento larval y generación de frass?
    \item ¿Qué hace de la BSF una alternativa biotecnológica atractiva para gestión de residuos?
\end{itemize}

\subsection*{Variables críticas del proceso}
\begin{itemize}[leftmargin=*]
    \item ¿Qué variables ambientales afectan más la bioconversión con BSF?
    \item ¿Por qué temperatura, humedad, pH y composición del sustrato son críticas?
    \item ¿Qué pasa si estas variables salen del rango adecuado?
    \item ¿Por qué ciertas condiciones favorecen la actividad larvaria?
    \item ¿Qué relación hay entre condiciones del entorno y calidad del proceso?
\end{itemize}

\subsection*{Emisiones y dinámica biológica}
\begin{itemize}[leftmargin=*]
    \item ¿En qué momentos del proceso esperas mayores emisiones de CO$_2$ o CH$_4$?
    \item ¿Las emisiones provienen de la larva, del sustrato o de la interacción entre ambos?
    \item ¿Cómo separarías el efecto biológico de la BSF del efecto propio del residuo?
    \item ¿Por qué conocer la dinámica biológica mejora la interpretación de las emisiones?
    \item ¿Qué variables biológicas mínimas debes entender para que tu modelo tenga sentido?
\end{itemize}

\subsection*{Sustrato, alimentación y rendimiento}
\begin{itemize}[leftmargin=*]
    \item ¿Todos los residuos orgánicos son igualmente aptos para BSF?
    \item ¿Cómo influye la composición del sustrato en el crecimiento larval?
    \item ¿Cómo afecta el sustrato la eficiencia de bioconversión y las emisiones?
    \item ¿Qué riesgos hay si el sustrato tiene exceso de humedad, acidez o heterogeneidad?
    \item ¿Cómo impacta la calidad del sustrato la calidad del frass y la biomasa?
\end{itemize}

\subsection*{Indicadores biológicos y límite del proyecto}
\begin{itemize}[leftmargin=*]
    \item ¿Qué indicadores biológicos o zootécnicos considerarías para decir que el proceso va bien?
    \item ¿Cómo relacionas crecimiento larval, reducción de residuos y emisiones?
    \item ¿Cuál sería una señal biológica temprana de que el sistema entra en estado crítico?
    \item ¿Tu tesis estudia la cría de BSF o la medición del proceso con BSF?
    \item ¿Qué aspectos zootécnicos quedan fuera del alcance de tu propuesta?
\end{itemize}

\section{Preguntas transversales de candidatura}

\begin{itemize}[leftmargin=*]
    \item ¿Cuál es exactamente el aporte doctoral de esta tesis?
    \item ¿Qué brecha de conocimiento cierra tu propuesta?
    \item ¿Cuál es la hipótesis central del trabajo?
    \item ¿Qué resultado negativo sería igualmente valioso científicamente?
    \item ¿Cuál es la principal originalidad metodológica?
    \item ¿Qué hallazgo esperas que sea publicable?
    \item ¿Qué componente de la tesis podría mantenerse aunque otro falle?
    \item ¿Cómo defenderías el proyecto frente a un jurado que diga que es demasiado ambicioso?
    \item ¿Cómo defenderías el proyecto frente a un jurado que diga que aún está poco delimitado?
    \item ¿Qué decisiones ya están definidas y cuáles siguen abiertas?
\end{itemize}

\section*{Sugerencia de estudio}
Preparar para cada pregunta una respuesta en tres niveles:
\begin{itemize}[leftmargin=*]
    \item respuesta corta de 30 segundos;
    \item respuesta estándar de 1 a 2 minutos;
    \item respuesta ampliada de 3 minutos con ejemplo.
\end{itemize}

\end{document}
